\section{Label sets}\label{app:sets}
Table~\ref{tab:classmap} gives the full class map.
Class names list the decay products, with $q$ for an up or down quark, $g$ for a gluon, $\tau_h$ for a tau
lepton that decays to hadrons and $\nu$ for a neutrino; a prong is one collimated visible decay product, so neutrinos leave none.
Each JetClass-II class belongs to one class of every coarser label set, so the \Kmid{} classes are unions of
JetClass-II classes and the \Kcoarse{} classes unions of the \Kmid.

\input{tables/table_classmap}

\subsection{Flavor-boundary groupings}
Three further groupings into \Kcoarse{} classes were pretrained with \RandFlavRuns{} runs each, to ask whether
one boundary between flavors accounts for the difference between the \Kcoarse-class set and the random
groupings.
The first assigns whole families of decays that differ only in flavor to groups at random, so that, unlike the
\Kcoarse-class set, its groups mix decays of different structure, and no boundary between two groups separates
two flavors of one family; the second
moves the four-quark decays that contain a bottom quark into another group; and the third moves as many
four-quark decays chosen at random.
On the pretraining classes a linear probe ranks them between the \Kcoarse-class model and the random groupings:
the grouping without a flavor boundary lies above the \Kcoarse-class model, the two with a boundary lie above
that grouping, and the random groupings lie above all three, each step being significant (table~\ref{tab:flavor}).
The groupings with a boundary differ from the one without in more than that boundary, because other decays had
to move to keep the class sizes, so this result cannot isolate the boundary and we draw no conclusion from it in
the main text.

\input{tables/table_flavor}

\section{Training details}\label{app:training}
The JetClass-II jets are clustered with the anti-$k_{\mathrm T}$ algorithm~\cite{antikt} from the particles of a
simulated CMS-like detector, with \LitJetClassTwoPileupMean{} pileup collisions on average.
We keep jets with $p_{\mathrm T}$ between \DesignPtMin{} and \DesignPtMax~GeV and a soft-drop mass, computed at
$z_{\mathrm{cut}} = \LitCmsSdZcut$ and $\beta = \LitCmsSdBeta$, between \DesignMsdMin{} and \DesignMsdMax~GeV.
The generator-level mass $m_{\mathrm{gen}}$ is the soft-drop mass of the jet of generator-level particles matched
to the reconstructed jet; jets without a match enter neither the regression term nor table~\ref{tab:mass}.
The quark--gluon jets are smaller, with a radius of \LitQgJetRadius, and have no detector simulation.
Each model sees \DesignExamplesSeen{} training examples and is evaluated at the epoch of its highest validation
accuracy among the last ten.
Pretraining follows Sophon's settings: \DesignEpochs{} epochs of \DesignSamplesPerEpoch{} examples in batches of
\DesignBatchSize, a learning rate of \DesignStartLr{} that is constant for the first \DesignLrFlatPercent{} of
the epochs and then decays to \DesignLrEndPercent{} of its value, and the particle momenta of every jet rescaled
to a common jet $p_{\mathrm T}$.
The fine-tuned network has \MlstNParamsJc{} parameters on JetClass.
On JetClass-II the linear probes are trained on \VProbeNTrain{} jets, at most \VProbePerClassMax{} per class, and
tested on \JcBalNTest, at most \JcBalTestPerClass{} per class, with the \JcBalNShort{} rarest classes contributing
all the jets they have; on JetClass they are trained on \LpNTrainMaxJci{} jets and on the two binary benchmarks
on the full training sets.
Fine-tuning on JetClass-II and JetClass uses a learning rate of \MlstFtLr{} for the pretrained weights and one
\MlstFtHeadMult{} times larger for the new output layer, a weight decay of \MlstFtWeightDecay{} and batches of
\MlstFtBatch{} jets, for \FtStepsESix{} steps at \FtSizeESix{} jets, and keeps the weights of the best
validation check; the learning rates of the other training sizes and of the network trained from scratch
were chosen on validation from a grid, and the full scan is in the data release.
All pretraining runs together took \MlstComputeTotalGpuHours{} GPU hours.

The binary probes on pairs of classes are tested on \JcBalPairTestPerSide{} jets of each class, and a curve of
figure~\ref{fig:pairs} is drawn only where at least \MlstMinBkgPass{} background jets pass.

For the resonance benchmark the transfer network has a hidden layer of \DesignHeadHidden{} units and two outputs.
Its hidden layer starts from that of the pretrained output layer and all of its weights are then trained, on
\XbbNTrainPerClass{} jets per class for $X\to b\bar b$ and \XbsNTrainPerClass{} for $X\to bs$, with
\MlstXbxHeadSeedsWord{} seeds per pretraining run, which are averaged.
The network trained from scratch for each signal is the Sophon network on the same inputs and jets, with
\MlstXbxScratchSeedsWord{} seeds.
The QCD jets come from files that no pretraining run read, and signal and background are counted in the same
window of soft-drop mass.
The significance of a difference on this benchmark uses a binomial estimate of the test-sample term, because
the per-jet scores were not kept.

\section{Inputs and missing inputs}\label{app:inputs}
Table~\ref{tab:inputs} lists the inputs and which of them the two binary benchmarks provide.
For the linear probes an absent input is filled with zero, with the training mean or with random draws of real
JetClass-II particles; the choice is made once per task on validation, for every model, and is
\LpSubstitutionTop{} for top tagging and \LpSubstitutionQg{} for quark--gluon discrimination.
\input{tables/table_inputs}

\section{Further results}\label{app:results}
Tables~\ref{tab:auc} and~\ref{tab:herwig} give the AUC of every model of tables~\ref{tab:probes}
and~\ref{tab:finetuned} and the quark--gluon results on jets from the Herwig generator~\cite{qgherwig}.

Tables~\ref{tab:head} to~\ref{tab:resonance} hold the numbers that the main text quotes from no other table:
the pretrained output layer, the linear probes' AUC by class group and for pairs of classes, the fine-tuned
accuracy at smaller training sets, and the resonance benchmark.

\input{tables/table_auc}
\input{tables/table_herwig}
\input{tables/table_head}
\input{tables/table_pairs}
\input{tables/table_size}
\input{tables/table_resonance}

\subsection{MLP probes}
An MLP probe with one hidden layer of \LpMlpHidden{} units, trained on the same frozen features as the linear
probes, tests whether a conclusion depends on the probe being linear (table~\ref{tab:mlp}).
It changes two of them.
When the \Kfull{} classes are merged into the \Kcoarse{} groups of the coarsest label set and the probe is asked
only for the group, linear probes show no significant difference between the label sets, whereas the MLP probe
finds the features of the \Kcoarse-class model significantly worse than those of the \Kfull-class model, on the
task its own labels describe (table~\ref{tab:groups}).
On top tagging the MLP probe separates one configuration that the linear probe does not, the \Kcoarse-class
model with mass regression from the \Kfull-class model.
The mass-regression nulls of the main text hold under the MLP probe.

\input{tables/table_mlp}
\input{tables/table_groups}

\section{Comparison with published taggers}\label{app:published}
Table~\ref{tab:published} places our fine-tuned and from-scratch networks beside published results on the three
public tasks.
The linear energy-flow-polynomial (EFP) classifier is linear in a set of jet observables; PFN, ParticleNet,
ParT, MIParT, L-GATr and OmniLearned are neural networks on the particles, the last four transformers.
On quark--gluon discrimination ours are close to the published Particle Transformer.
On top tagging our network trained from scratch is below the published one trained from scratch, and our
fine-tuned network is above that and below the published Particle Transformer fine-tuned from JetClass.
That one was pretrained on the kinematic inputs only; our \Kfull-class network pretrained that way, one run,
reaches the same rejection after fine-tuning as with the absent inputs set to zero (last row of the table).

\input{tables/table_published}

\section{Further figures}\label{app:figures}
Figure~\ref{fig:recall} shows which classes lose recall under each coarser label set, figure~\ref{fig:otherpairs}
the class pairs that figure~\ref{fig:pairs} omits, figure~\ref{fig:probesize} the linear probes at smaller training
sets, figure~\ref{fig:masspt} the regressed mass against jet $p_{\mathrm T}$ and mass, and
figure~\ref{fig:validation} the validation curves of pretraining.
\vfig{figA1}{Recall of each JetClass-II class under a linear probe on the frozen features of the (a)
\Kres-, (b) \Kmid- and (c) \Kcoarse-class model, against its recall for the \Kfull-class model. Points are means
over runs, with bars combining run spread and test-sample uncertainty; colors and shapes give what the panel's label
set merges the class with. Open points lie within three standard deviations of equal recall.}{fig:recall}
\vfig{figA2}{As figure~\ref{fig:pairs}, for the four other class pairs: (a) $X\to YY\to b\bar bqq$ against
$c\bar cqq$, and QCD jets (b) with $b\bar b$ against $c\bar c$, (c) with $b$ against light quarks and gluons only,
(d) with $c$ and no $b$ against light quarks and gluons only. Every random grouping merges each QCD pair, so (b) to
(d) have one random-grouping curve.}{fig:otherpairs}
\vfig{figA3}{Linear probes on the frozen features against the number of training jets: accuracy on (a)
JetClass-II and (b) JetClass, and background rejection Rej$_{30\%}$ on (c) top tagging and (d) quark--gluon
discrimination. Lines are the mean over runs and bands the standard deviation.}{fig:probesize}
\vfig{figA4}{Jet mass estimated by the models pretrained with mass regression, beside the soft-drop mass:
median response (a, b) and resolution (c, d) against jet $p_{\mathrm T}$ and against generator-level mass. Lower
panels give the distance of the median from one and the ratio of the resolution to that of the soft-drop mass. Bars
combine run spread and test-sample uncertainty.}{fig:masspt}
\vfig{figA5}{Validation accuracy of the pretraining runs against epoch, each model on its own label set,
so levels compare only between models with the same number of classes. Lines are the mean over runs and bands the
standard deviation.}{fig:validation}
